% =====================================================================
% Macros de apoio aos slides
% =====================================================================

% Marcador visível de conteúdo pendente. Some ao compilar em modo final
% ("make final"), para não ir pendência nenhuma na versão de entrega.
\newif\ifrascunho
\rascunhotrue
\ifdefined\modofinal\rascunhofalse\fi

\newcommand{\pendente}[1]{%
  \ifrascunho\textcolor{peeaccent}{\footnotesize\textbf{[TODO: #1]}}\fi}

% Destaque semântico.
\newcommand{\dest}[1]{\textcolor{peedark}{\textbf{#1}}}
\newcommand{\alerta}[1]{\textcolor{peeaccent}{\textbf{#1}}}

% Figura ocupando o slide inteiro, com legenda opcional no rodapé.
% Uso: \framefigura{caminho/figura}{Legenda}
\newcommand{\framefigura}[2]{%
  \begin{frame}[plain]
    \centering
    \includegraphics[width=0.92\paperwidth,height=0.82\paperheight,keepaspectratio]{#1}\\[0.3em]
    {\scriptsize\color{peegray}#2}
  \end{frame}}

% Slide de dois blocos lado a lado.
% Uso: \duascolunas{esquerda}{direita}
\newcommand{\duascolunas}[2]{%
  \begin{columns}[T,onlytextwidth]
    \begin{column}{0.48\textwidth}#1\end{column}
    \begin{column}{0.48\textwidth}#2\end{column}
  \end{columns}}

% Abreviações recorrentes do domínio.
\newcommand{\TSAD}{\mbox{TSAD}}
\newcommand{\DQM}{\mbox{DQM}}
\newcommand{\LAr}{\mbox{LAr}}
